\chapter{Discussion}
\label{chap:discussion}

The proposed correction has a clear mathematical purpose: it represents
sampling uncertainty in the MI variance estimate before using that estimate
in a two-sample test. The experiments show that the purpose is narrower than
its possible practical effect. A better description of the denominator does
not by itself repair every finite-sample error in the numerator, in the
first-order approximation, or in the availability of a valid statistic.

\section{Answers to the research questions}

The first research question was whether an MI-specific effective degree of
freedom can be derived. Chapter~\ref{chap:expanded} does so by differentiating
the variance of pointwise MI along one-cell probability perturbations. The
derivative includes changes to row and column margins, not just to the
selected cell. Its sampling variance supplies the component degrees of
freedom used in Satterthwaite moment matching. These derivative and
moment-matching steps define a reproducible statistic. The final Student
reference remains a working approximation, not an exact
finite-sample theorem for MI\@.

The second question concerned false positives, detection and validity.
Figure~\ref{fig:main-8x8-uniform} contains a fully valid small-sample null
case where Expanded Welch moves an excessive Wald false-positive rate toward
0.05. Figure~\ref{fig:main-2x2-uniform} contains a larger-sample null case
where the same adjustment moves a below-nominal Wald rate farther from
0.05. In the sparsest panels, a low unconditional rejection rate may largely
reflect invalid outputs. Thus improved calibration at selected null points
does not imply that the expanded procedure performs better overall. Under
alternatives, its rejection rate is generally lower; these comparisons are
not power at a matched achieved size.

The third question asked whether population features alter this judgement.
The main landscape varies alphabet size and margins explicitly. The focused
comparisons show that baseline MI, changed-cell location, sample allocation
and population constructor can change the finite-sample behaviour even when
the target MI difference is held fixed. The construction comparison is
particularly instructive: equal margins and MI do not determine the table's
complete sampling distribution. A scalar target is useful for organising
experiments, but cannot guarantee identical difficulty across distinct
population tables. The \(8\times8\) imbalance check adds a related warning:
when \(d\) is large and one sample remains small, the leading bias correction
itself can be under substantial finite-sample stress. Changing only the
reference distribution cannot repair that shared numerator problem.

The fourth question concerned sample size and computation. For the tested
first-block positive-MI populations with true MI 0.02, both null rates are
near 0.05 at \(n=50{,}000\). The corresponding very-small-MI fixed
populations remain noticeably conservative at that sample size. Expanded
Welch has the same linear-in-cell-count order as Wald but costs about 1.75
times as much per scalar call in the measured implementation. The extra
expense is modest compared with a full resampling procedure but not cost-free,
particularly for many repeated comparisons.

\section{What the correction can and cannot change}

Both methods compute the same \(\Deltahat\), \(\SEhat\) and \(T\). With
valid positive degrees of freedom, Expanded Welch has a heavier-tailed
reference and a two-sided p-value no smaller than Wald's for the same table
pair. This gives a direct explanation for many plots: the correction lowers
rejection whether the Wald rate is too high or already too low. It is not an
adaptive calibration method that can move the threshold in both directions.
It also does not alter the leading bias correction or re-estimate the
numerator's finite-sample distribution.

This direction contrasts with the regression settings that motivate several
Satterthwaite precedents. There, downward-biased uncertainty estimates or
undercoverage call for a correction that makes rejection harder. Near MI
independence, however, \(\E\widehat V\approx V+d/n\), so estimated uncertainty
is already shifted upward; adding a heavier-tailed reference compounds the
conservatism instead of correcting it.

The estimated numerator and denominator come from the same count tables.
Independence of the \(P\) and \(Q\) \emph{samples} allows their variance
contributions to be combined, but it does not make the numerator independent
of its estimated standard error. Moment matching of the denominator alone
therefore does not prove that \(T\) has a Student distribution. When the
first-order MI approximation itself is poor, changing only the reference
quantile has limited scope. The independent-pilot diagnostic in
Appendix~\ref{app:mechanism} shows that replacing the random denominator by
the actual finite-sample standard deviation can restore near-nominal rejection
in a representative conservative regime. It also shows that the population
first-order standard deviation is too small there. The result supports a
combined explanation involving denominator estimation and omitted quadratic
numerator variation; it does not isolate mean denominator bias from
numerator--denominator dependence or establish one universal mechanism.
The component degree-of-freedom formula also has its own regularity
condition: a positive first-order MI variance does not force its variance
sensitivity to be positive. At such stationary points, the first-order
moment match cannot describe the finite-sample variance estimate without
higher-order terms.

Near independence, the local calculation gives \(V\approx2I\),
\(\tau^2\approx4V\), and hence a population component degree of freedom near
\(nI\). This explains why the Expanded Welch reference can become extremely
heavy-tailed at low baseline MI. The same local Gaussian calculation shows
that quadratic terms of dimension \(d\) affect both the mean and variance of
\(\widehat V\). First-order moment matching therefore requires \(nV\) to be
large relative to \(d\), in addition to adequate expected counts. The
follow-up comparison with simulated moments confirms that the population,
plug-in and finite-sample component degrees of freedom can differ materially
at \(n=100\), then move closer by \(n=1000\).

This transition is consistent with the broader second-order delta-method
description of \citet{marinescu2025}: away from independence the MI estimation
error contains a normal contribution together with quadratic, chi-squared-type
contributions, whereas the linear contribution vanishes at independence. That
result does not supply a finite-sample correction for the present two-sample
weak null, but it provides a plausible explanation for why a purely
first-order Student adjustment can remain conservative at small positive MI.

At exact independence, \(I=0\) and the population first-order MI variance
vanishes. This does not mean finite empirical MI estimates are always zero.
It means their leading stochastic variation is quadratic rather than linear
in the probability errors. The standard one-table likelihood-ratio statistic
uses that second-order behaviour. Neither setting a reference \(Q\) to an
independent distribution nor squaring the current \(T\) makes Expanded Welch
the classical independence test. Appendix~\ref{app:independence-boundary}
gives the mathematical distinction.

\section{Interpretation of invalid samples}

The validity diagnostics matter most where expected counts are very small.
Wald needs a usable total estimated variance. The expanded implementation
additionally needs usable component sensitivity variances and effective
degrees of freedom. An observed table can have a positive total variance
while one component calculation is unusable. The final study does not silently
replace these cases with a different test. Instead, it reports invalidity
and counts such a replicate as a non-rejection in the primary unconditional
rate.

That rate answers an operational question: how often does the complete
procedure produce a valid rejection? It should not be mistaken for the
conditional rejection probability among valid tables. Nor does a failure to
return a p-value constitute a true negative. The saved method-valid and
common-valid subsets allow these distinctions to be inspected. In
particular, the very large \(8\times8\), \(n=5\) alternative gap in
Chapter~\ref{chap:results} cannot be explained as a pure change of reference
threshold because the methods are valid on different fractions of tables.
The exact paired records are preferable to comparing two method-specific
conditional rates with different denominators.

\section{Practical conclusion and limits of the evidence}

For the tested task, Normal Wald remains the simpler general baseline.
Expanded Welch can be useful as a conservative sensitivity analysis in a
regime where its validity is high and Wald's null rejection is demonstrably
liberal. It should not be installed as a universal replacement merely
because one null table shows a smaller false-positive rate. In applications,
the population truth is unknown, so the most relevant calibration study is
one matched to the anticipated alphabet, sample size and plausible margins.
Selecting between the methods after inspecting a single p-value is not a
validated selection rule.

Several limitations bound this conclusion. The population families are
deliberately constructed fixed examples, not a probability sample of
all joint distributions. Twenty thousand replicates make a rate at a fixed
population relatively precise; they do not make the 534 selected population
pairs representative of every application. The tables are at most
\(8\times8\) square and \(3\times5\) rectangular, with the stated finite
sample ladders and only a two-sided nominal level of 0.05. The main
construction emphasises a local four-cell association; spread, rare-cell and
log-linear checks probe but do not exhaust other possibilities. The study
does not validate paired observations, continuous estimators, a real-data
workflow or growing alphabets.

The protocol was frozen before the final simulation, after exploratory
development of the method and design. It is therefore a reproducible final
run, not an externally preregistered untouched hypothesis. Configurations
displayed in several sections are the same sampled run. Repeating their
figure appearances does not provide independent confirmation. The recorded
Monte Carlo intervals also address repeated sampling of each fixed pair, not
the broader uncertainty of generalising to an untested population.

The reported alternative rejection rates use the common nominal 0.05
threshold. Size-adjusting each method separately would answer a different
question: how their rankings change when a simulation-specific critical value
is known. Such a curve would not provide a usable calibration rule for an
unknown real population and was not part of the frozen comparison. The thesis
therefore keeps calibration and nominal-threshold detection visible side by
side and does not describe the alternative curves as power at matched size.

Future work should address the joint finite-sample behaviour of the
numerator and denominator or use a calibrated reference tailored to sparse
tables. The transition near independence is a particularly distinct problem
because the first-order variance tends to zero. A likelihood-ratio or
resampling alternative must be assessed under the equal-MI \emph{weak null},
not imported from a test of full distributional equality or one-table
independence without checking its target. These are directions for new
validation, not remedies established by the present results.
